\begin{center}
	\LARGE
	Bag en kage
\end{center}

I skal i denne aflevering bage en kage. Det må være i netværksgrupper eller grupper af eget valg. Mængden af ingredienser skal I finde ved at løse en række opgaver. Opskriften er på en gulerodskage med cream cheese frosting tyvstjålet fra en ikke nærmere bestemt, men ganske udbredt dansk madblogger.

\textbf{Ingredienser til kagen}
\begin{itemize}
	\item[$\cdot$] \underline{\phantom{150}} g stuetempereret smør.
	\item[$\cdot$] \underline{\phantom{100}} g sukker.
	\item[$\cdot$] \underline{\phantom{100}} g brun farin.
	\item[$\cdot$] \underline{\phantom{4}} æg.
	\item[$\cdot$] \underline{\phantom{200}} g hvedemel.
	\item[$\cdot$] \underline{\phantom{2}} tsk bagepulver.
	\item[$\cdot$] \underline{\phantom{1}} tsk kanel.
	\item[$\cdot$] \underline{\phantom{0.5}} tsk salt.
	\item[$\cdot$] \underline{\phantom{275}} g skraldede og groftrevne gulerødder.
\end{itemize}

\textbf{Ingredienser til topping}
\begin{itemize}
	\item[$\cdot$] \underline{\phantom{125}} g stuetempereret smør.
	\item[$\cdot$] \underline{\phantom{1}} tsk vaniljesukker 
	\item[$\cdot$] \underline{\phantom{125}} g flormelis.
	\item[$\cdot$] \underline{\phantom{125}} g naturel flødeost.
	\item[$\cdot$] Saft og skal fra \underline{\phantom{1}} økologisk citron.
\end{itemize}

\textbf{Fremgangsmåde til kage}
Pisk smør, sukker og brun farin luftigt. Pisk derefter ét æg i ad gangen, så blandingen ikke skiller. 

Bland de tørre ingredienser godt sammen i en anden skål og vend herefter de tørre ingredienser i dejen. Vend til sidst den revne gulerod i dejen og bag kagen i en smurt springform på 20cm diameter ved 170$^\circ$C varmluft i ca. 45 minutter. 

\textbf{Fremgangsmåde til frosting}
Pisk smør, sukker og vaniljesukker sammen til det er luftigt. Bland citronsaft, citronskal og flødeost i blandingen og smør den på kagen, når kagen er helt kold. 

\newpage

\begin{opgavetekst}{Opgave 1}
\end{opgavetekst}	
\begin{delopgave}{1}
	\begin{align*}
		\frac{601}{4} + \frac{52}{2} - \frac{210}{8}
	\end{align*}
\end{delopgave}
%
%
\begin{opgavetekst}{Opgave 2}
	Et polynomium $p$ er givet ved
	\begin{align*}
		p(x) = ax^2 - 500x + 600.
	\end{align*}
\end{opgavetekst}
\begin{delopgave}{1}
	Bestem tallet $a$, så $p$ har rødder i punkterne $x = 2$ og $x = 3$.  
\end{delopgave}
\begin{meretekst}
	Et andet polynomium $q$ er givet ved
	\begin{align*}
		q(x) = x^2 + 99x + c
	\end{align*}
\end{meretekst}
\begin{delopgave}{2}
	Bestem tallet $c$, så parablerne for $p$ og $q$ skærer hinanden når $x = 1$. 
\end{delopgave}
%
%
\begin{opgavetekst}{Opgave 3}
	To linjer $l$ og $m$ er givet ved ligningerne
	\begin{align*}
		&l: \ -4x + 2y = -20 \\
		&m: \ 5x + y = 18
	\end{align*}
\end{opgavetekst}
\begin{delopgave}{1}
	Bestem $x$-koordinaten til skæringspunktet mellem linjerne $l$ og $m$. 
\end{delopgave}
\begin{delopgave}{2}
	Bestem $y$-koordinaten for $l$ når $x = 105$. 
\end{delopgave}
%
%
\begin{opgavetekst}{Opgave 4}
	En funktion $f$ er givet ved
	\begin{align*}
		f(x) = \ln(2x-3).
	\end{align*}
\end{opgavetekst}
\begin{delopgave}{1}
	Bestem $f^{-1}(0)$
\end{delopgave}
%
%
\begin{opgavetekst}{Opgave 5}
	En eksponentialfunktion $g$ er givet ved
	\begin{align*}
		g(x) = 7\cdot 1.01^x
	\end{align*}
\end{opgavetekst}
\begin{delopgave}{1}
	Afgør, hvor mange $\%$ $g$ vokser med, hvis $x$ øges med 1. 
\end{delopgave}
%
%
\begin{opgavetekst}{Opgave 6}
	Et polynomium $q$ er givet ved
	\begin{align*}
		q(x) = ax^2 + 5x + 12.5
	\end{align*}
\end{opgavetekst}
\begin{delopgave}{1}
	Bestem tallet $a$, så $q$ har netop én rod. 
\end{delopgave}
%
%
\begin{opgavetekst}{Opgave 7}
	En eksponentialfunktion $h$ er givet ved
	\begin{align*}
		h(x) = 2\cdot 0.99748^x
	\end{align*}
\end{opgavetekst}
\begin{delopgave}{1}
	Bestem halveringskonstanten for $h$.
\end{delopgave}
\begin{delopgave}{2}
	Løs ligningen $h(x) = 1.46$
\end{delopgave}
%
%
\begin{opgavetekst}{Opgave 8}
	En funktion $f$ har forskriften
	\begin{align*}
		f(x) = \sqrt{5x}.
	\end{align*}
\end{opgavetekst}
\begin{delopgave}{1}
	Hvor langt skal grafen for $f$ parallelforskydes til venstre for at gå gennem punktet $(4,5)$? 
\end{delopgave}
%
%
\begin{opgavetekst}{Opgave 9}
	I en bank får du en årlig rente på 2.8$\%$. Du indsætter $100\TS 000$ kr i banken.
\end{opgavetekst}
\begin{delopgave}{1}
	Hvor mange penge står på kontoen efter 8 år? Beløbet skal angives i 1000 kroner. 
\end{delopgave}
\begin{delopgave}{2}
	Efter hvor mange år vil der være $3 \TS 156 \TS 016$ kr på kontoen?
\end{delopgave}
%
%
\begin{opgavetekst}{Opgave 10}
	En stykvist defineret funktion $g$ er givet ved
	\begin{align*}
		g(x) = 
		\begin{cases}
			x^2 + 4x - 3,  &\textnormal{ hvis }x \leq 2,\\
			4x + b,  &\textnormal{ hvis }x > 2.
		\end{cases}
	\end{align*}
\end{opgavetekst}
\begin{delopgave}{1}
	Bestem tallet $b$, så grafen for $g$ er sammenhængende. 
\end{delopgave}